\documentclass[11pt,a4paper]{article}
\usepackage[margin=2.4cm]{geometry}
\usepackage{amsmath, amssymb, amsthm, mathtools}
\usepackage{xcolor}
\usepackage{microtype}
\usepackage{enumitem}
\usepackage{hyperref}
\hypersetup{colorlinks, linkcolor=blue!50!black, citecolor=blue!50!black,
            urlcolor=blue!60!black}

\newcommand{\eml}{\mathrm{eml}}

\title{Why a single binary operator deserves attention\\[3pt]
       \large A speculative companion to the EML formalization}

\author{Bartosz Naskręcki\\[2pt]
  \small Adam Mickiewicz University, Poznań / Politechnika Warszawska\\
  \small \texttt{bartosz.naskrecki@amu.edu.pl}}

\date{}

\begin{document}
\maketitle

\begin{abstract}
This is a deliberately speculative companion to the technical
formalization paper
(\texttt{notes/extended\_paper/eml\_formalization\_paper.pdf}). The
formalization establishes that one binary operator,
$\eml(x, y) := \exp(x) - \log(y)$, paired with the constant $1$,
generates the $36$-element catalogue of elementary functions used in
typical scientific calculators
\cite{odrzywolek2026eml}. The purpose of \emph{this} document is
to argue that the existence of such a continuous Sheffer operator is
not a curiosity of the elementary-function algebra alone: similar
``one binary plus one terminal'' generative reductions arise across
several distinct areas of physics --- some explicitly, some
implicitly --- and the EML observation can be used to clarify what
structural fact is being relied on in each case.

We make no claim of formal isomorphism; the analogies offered here
are structural prompts, not theorems. The reader is invited to take
them as a programme of \emph{questions} for the relevant
sub-communities, rather than as established connections.
\end{abstract}

\tableofcontents

% =====================================================================
\section{Scope and disclaimer}
\label{sec:scope}
% =====================================================================

This note is exploratory. The technical paper
\cite{naskrecki2026formalization} establishes a precise mathematical
fact (a Lean 4 + Mathlib formalization of Odrzywołek's completeness
theorem); the present document is concerned with whether that fact
points to a broader pattern. We separate the two so that the
technical paper can stand on its formal-methods merits without being
diluted by speculative analogies.

The four sections that follow each pose the question \emph{``in
domain $X$, is there a structural reduction analogous to the EML
Sheffer operator?''} and offer a sketch of what an affirmative
answer would look like. Sections \ref{sec:hamiltonian} (Hamiltonian
normal forms), \ref{sec:wess-zumino} (Wess--Zumino covering
identities), \ref{sec:cft} (conformal-bootstrap reductions), and
\ref{sec:analog-computing} (single-instruction analog hardware) are
each a few pages of motivated speculation. Section
\ref{sec:open-questions} collects the open problems that fall out
of the analogies.

% =====================================================================
\section{Hamiltonian normal forms and the Sheffer reduction}
\label{sec:hamiltonian}
% =====================================================================

[TO BE WRITTEN — develops the analogy between (a) the EML reduction
of the 36-primitive elementary toolbox to a single binary operator
and (b) the reduction of a generic Hamiltonian system to its
Birkhoff--Gustavson normal form, where a complicated dynamical
system collapses to a finite expression in a small canonical
generator set. The point is: in both cases the existence of a small
generative basis is non-trivial and has computational value.

Open question: does the Birkhoff--Gustavson normal form admit a
``Sheffer reduction'' --- a single canonical-transform kernel from
which all integrable normal forms can be recovered, paired with a
finite seed? If yes, what is the kernel?]

% =====================================================================
\section{Wess--Zumino-style covering identities}
\label{sec:wess-zumino}
% =====================================================================

[TO BE WRITTEN — discusses the 4d Wess--Zumino superalgebra's
generator reduction (the entire algebra is built from a single odd
spinor charge $Q_\alpha$ and its conjugate, plus the bracket; analogue
to ``one binary plus one terminal''). Comments on whether
Sheffer-style reductions exist for the bosonic generators of the
super-Poincaré algebra.

Open question: is the supercharge's generative role essential, or is
it a notational convenience? In the EML case we can prove (via the
formalization) that no smaller generator set works for the full
36-primitive catalogue; what is the analogous statement in the
supercharge setting?]

% =====================================================================
\section{Gauge-fixing and conformal-bootstrap reductions}
\label{sec:cft}
% =====================================================================

[TO BE WRITTEN — Belavin--Polyakov--Zamolodchikov's conformal
bootstrap reduces the infinite-dimensional space of CFT correlators
to a small set of structure constants and conformal blocks; the
analogue is a Sheffer-style reduction in the algebra of correlators.

Open question: in the bootstrap algebra, is the OPE the ``single
binary operator'' analogue, with the identity operator playing the
role of the constant 1? If so, what is the analogue of the
36-primitive catalogue --- the list of correlators that the
bootstrap is required to reproduce?]

% =====================================================================
\section{Single-instruction analog computing}
\label{sec:analog-computing}
% =====================================================================

[TO BE WRITTEN — discusses FPGAs, photonic integrated circuits, and
the practical implication of the EML Sheffer: a ``single-instruction
math co-processor'' that evaluates only $\eml(a, b)$ in hardware. The
formal proof of completeness implies that such a co-processor is
universal for the 36-primitive class, modulo the boundary points;
this is a substantially stronger statement than the empirical
verification provides.

Open question: what is the silicon-area / latency / energy tradeoff
of a single-EML-instruction analog co-processor, compared with the
standard floating-point unit that implements all 36 primitives
separately? At what circuit size does the Sheffer reduction become
economically interesting?]

% =====================================================================
\section{Open questions}
\label{sec:open-questions}
% =====================================================================

The four analogies above each pose a concrete open question
(highlighted at the end of the corresponding section). Collected
here for ease of reference:

\begin{enumerate}[topsep=2pt, itemsep=4pt]
  \item \textbf{Birkhoff--Gustavson kernel.} Does the
    Birkhoff--Gustavson normal form admit a single canonical-transform
    kernel from which all integrable normal forms can be recovered,
    paired with a finite seed?

  \item \textbf{Supercharge minimality.} In the 4d
    Wess--Zumino algebra, can the role of the spinor supercharge as
    ``the single generator'' be made into a formal Sheffer-style
    minimality theorem?

  \item \textbf{OPE as Sheffer operator.} In the conformal bootstrap,
    can the operator product expansion be cast as a Sheffer
    operator on the algebra of correlators, with the identity
    operator as constant?

  \item \textbf{Single-EML hardware.} What is the practical viability
    of a single-instruction EML analog co-processor for elementary
    function evaluation?
\end{enumerate}

% =====================================================================
\section{What the formalization tells us}
\label{sec:what-formalization-says}
% =====================================================================

[TO BE WRITTEN — closing observation. The formalization makes two
contributions to the broader programme: (a) it establishes that
\emph{the EML Sheffer reduction is a theorem, not a conjecture}, so
analogues in physics can be discussed without the loophole "but is
it really minimal?"; (b) the witness-family form (one term per
region of the input domain) is precisely the form that a
single-instruction hardware unit would naturally implement, since
hardware does case analysis on input range as a matter of course.

The technical paper restricts itself to the formalization. This note
exists to argue that the formalization's value extends beyond
formal methods: a structural fact about elementary functions is
also a structural prompt for several adjacent fields, and the prompt
is now mathematically settled.]

\medskip

\begin{thebibliography}{9}

\bibitem{odrzywolek2026eml}
A.~Odrzywołek,
\emph{All elementary functions from a single binary operator},
arXiv:2603.21852, 2026.
\url{https://arxiv.org/abs/2603.21852}

\bibitem{naskrecki2026formalization}
B.~Naskręcki,
\emph{Formalizing the EML completeness theorem in Lean 4: an
extended account of the construction and of the multi-agent
workflow that produced it},
companion paper
\texttt{notes/extended\_paper/eml\_formalization\_paper.pdf} in the
\href{https://github.com/nasqret/eml-formalization}{eml-formalization}
repository, 2026.

\end{thebibliography}

\end{document}
